\documentclass[11pt,a4paper]{article}

%% PACHETE MATEMATICE
\usepackage{amsmath,amssymb,amsfonts}
\usepackage{amsthm}
\usepackage{mathpazo}
\usepackage{eulervm}
\usepackage{enumerate}
\usepackage{color}
\usepackage{graphicx}
\usepackage[all]{xy}
\usepackage{xcolor}
\usepackage{framed}
\usepackage[unicode]{hyperref}
\setcounter{MaxMatrixCols}{20}
\usepackage{multicol}
% \usepackage[romanian]{babel}


%% ASPECT
\linespread{1.05}
\usepackage[T1]{fontenc}
\usepackage[utf8x]{inputenc}
\usepackage[margin=2cm]{geometry}
% \usepackage[color]{showkeys}
\usepackage{anyfontsize}
\usepackage{titlesec}
\titleformat*{\section}{\large\bfseries}

\usepackage{fancyhdr}
\pagestyle{fancy}
\rhead{\small \color{gray} Notițe de Adrian Manea}
\lhead{\small \color{gray} ETTI-AM2, 2017---2018, M. Joița \& A. Niță}


\usepackage[autostyle = false, style = german]{csquotes}
\usepackage[romanian]{babel}
\newcommand\qq{\enquote}

%% COMENZILE MELE
\renewcommand*{\proofname}{Dem.:}
\newcommand\bb{\mathbb}
\newcommand\dr{\mathrm}
\newcommand\kal{\mathcal}
\newcommand\fk{\mathfrak}
\newcommand\op{\oplus}
\newcommand\ot{\otimes}
\newcommand\ds{\displaystyle}
\newcommand\id{\indent}
\newcommand\nid{\noindent}
\newcommand\seq{\subseteq}
\newcommand\sneq{\subsetneq}
\newcommand\speq{\supseteq}
\newcommand\spneq{\supsetneq}
\newcommand\rar{\rightarrow}
\newcommand\Rar{\Rightarrow}
\newcommand\xrar{\xrightarrow}
\newcommand\lar{\leftarrow}
\newcommand\Lar{\Leftarrow}
\newcommand\xlar{\xleftarrow}
\newcommand\lrar{\leftrightarrow}
\newcommand\Lrar{\Leftrightarrow}
\newcommand\ideal{\trianglelefteq}
\newcommand\ai{\text{ a.\^{i}. }}
\newcommand\sk{(\textit{Schi\c{t}\u{a}}) }
\newcommand\rhup{\rightharpoonup}
\newcommand\rhdn{\rightharpoondown}
\newcommand\lhup{\leftharpoonup}
\newcommand\lhdn{\leftharpoondown}
\newcommand\vid{\emptyset}
\newcommand\wh{\widehat}
\newcommand\ol{\overline}
\newcommand\NN{\mathbb{N}}
\newcommand\RR{\mathbb{R}}
\newcommand\ZZ{\mathbb{Z}}
\newcommand\QQ{\mathbb{Q}}
\newcommand\CC{\mathbb{C}}

%% STILURILE MELE
\newtheoremstyle{dotless}{}{}{\itshape}{}{\bfseries}{:}{ }{}

\theoremstyle{dotless}
\newtheorem{theorem}{Teorem\u{a}}[section]
\newtheorem{proposition}{Propozi\c{t}ie}[section]
\newtheorem{lemma}{Lem\u{a}}[section]
\newtheorem{corollary}{Corolar}[section]
\newtheorem{exercise}{Exerci\c{t}iu}

\newtheoremstyle{dotlessdr}{}{}{}{}{\bfseries}{:}{ }{}
\theoremstyle{dotlessdr}
\newtheorem{example}{Exemplu}[section]
\newtheorem{definition}{Defini\c{t}ie}[section]
\newtheorem{remark}{Observa\c{t}ie}[section]




\begin{document}
% \definecolor{labelkey}{RGB}{222,0,0}

\begin{center}
  {\large\textbf{Seminar 7 \\
      Ecuații cu derivate parțiale \\
      Metode de rezolvare}}
\end{center}

\vspace{2cm}

\section{Ecuații cu coeficienți variabili}

\indent\indent Ecuațiile cu derivate parțiale de ordinul al doilea, cu coeficienți variabili,
se rezolvă similar celor cu coeficienți constanți. Vom prezenta două exemple.

\textbf{Exemplu 1:} Să se aducă la forma canonică ecuația:
\[
  (1 + x^2)\frac{\partial^2 u}{\partial x^2} + %
  (1 + y^2)\frac{\partial^2 u}{\partial y^2} + %
  x \frac{\partial u}{\partial x} + y \frac{\partial u}{\partial y} = 0.
\]

\emph{Soluție:} Deoarece avem $ AC - B^2 = (x^2 + 1)(y^2 + 1) > 0 $, rezultă
că ecuația este de tip eliptic.

Ecuația caracteristică este:
\[
  (1 + x^2)dy^2 + (1 + y^2) dx^2 = 0,
\]
care înseamnă:
\[
  \sqrt{1 + x^2} dy = \pm i \sqrt{1 + y^2} dx.
\]
Rezultă că familiile de curbe caracteristice sînt:
\[
  \begin{cases}
    \ln(y + \sqrt{1 + y^2}) + i \ln(x + \sqrt{1 + x^2}) &= c_1 \\
    \ln(y + \sqrt{1 + y^2}) - i \ln(x + \sqrt{1 + x^2}) &= c_2
  \end{cases}
\]
În consecință, facem schimbarea de variabile:
\[
  \begin{cases}
    \tau &= \ln(y + \sqrt{1 + y^2}) \\
    \eta &= \ln(x + \sqrt{1 + x^2})
  \end{cases}
\]
Derivatele parțiale în funcție de noile variabile sînt:
\begin{align*}
  \frac{\partial u}{\partial x} &= \frac{\partial u}{\partial \eta} \cdot %
                                  \frac{1}{\sqrt{1 + x^2}} \\
  \frac{\partial^2 u}{\partial x^2} &= \frac{1}{1 + x^2} \cdot \frac{\partial^2 u}{\partial \eta^2}%
                                      - \frac{x}{(1 + x^2)^{\frac{3}{2}}} \cdot %
                                      \frac{\partial u}{\partial \eta} \\
  \frac{\partial u}{\partial y} &= \frac{1}{\sqrt{1 + y^2}} \cdot \frac{\partial u}{\partial \tau} \\
  \frac{\partial^2 u}{\partial y^2} &= \frac{1}{1 + y^2} \cdot \frac{\partial^2 u}{\partial \tau^2}%
                                      - \frac{y}{(1 + y^2)^{\frac{3}{2}}} \cdot %
                                      \frac{\partial u}{\partial \tau}.
\end{align*}

Rezultă că ecuația se reduce la forma canonică:
\[
  \frac{\partial^2 u}{\partial \eta^2} + \frac{\partial^2 u}{\partial \tau^2} = 0.
\]

\vspace{1cm}

\textbf{Exemplu 2:} Să se aducă la forma canonică și să se determine soluția
generală a ecuației:
\[
  x^2 \frac{\partial^2 u}{\partial x^2} - 2xy \frac{\partial^2 u}{\partial x\partial y} + %
  y^2 \frac{\partial^2 u}{\partial y^2} + x \frac{\partial u}{\partial x} + %
  y \frac{\partial u}{\partial y} = 0.
\]

\emph{Soluție:} Deoarece $ A = x^2, B = -xy, C = y^2 $, avem $ AC - B^2 = 0 $, deci
ecuația este de tip parabolic.

Din ecuația caracteristică obținem:
\[
  x^2 \cdot \Big(\frac{dy}{dx}\Big)^2 + 2xy \cdot \frac{dy}{dx} + y^2 = 0 \Rightarrow %
  \frac{dy}{dx} = -\frac{y}{x} \Rightarrow xy = c.
\]
Facem schimbarea de variabile:
\[
  \begin{cases}
    \tau &= xy \\
    \eta &= x
  \end{cases}
\]
și noile derivate parțiale sînt:
\begin{align*}
  \frac{\partial u}{\partial x} &= y\frac{\partial u}{\partial \tau} + %
                                  \frac{\partial u}{\partial \eta} \\
  \frac{\partial u}{\partial y} &= x \frac{\partial u}{\partial \tau} \\
  \frac{\partial^2 u}{\partial x^2} &= y^2 \frac{\partial^2 u}{\partial \tau^2} + %
                                      2y \frac{\partial^2 u}{\partial \tau\partial \eta} + %
                                      \frac{\partial^2 u}{\partial \eta^2} \\
  \frac{\partial^2 u}{\partial y^2} &= x^2 \frac{\partial^2 u}{\partial \tau^2} \\
  \frac{\partial^2 u}{\partial x\partial y} &= \frac{\partial u}{\partial \tau} + %
                                              xy \frac{\partial^2 u}{\partial \tau^2} + %
                                              x \frac{\partial^2 u}{\partial \tau \partial \eta}.
\end{align*}

Atunci ecuația devine:
\[
  \frac{\partial^2 u}{\partial \eta^2} + \frac{1}{\eta} \frac{\partial u}{\partial \eta} = 0.
\]
Ea se poate rescrie și rezolva astfel:
\[
  \frac{\partial}{\partial \eta} \Big( \eta \cdot \frac{\partial u}{\partial \eta} \Big) = 0 %
  \Rightarrow \eta \cdot \frac{\partial u}{\partial \eta} = f(\tau) \Rightarrow %
  \frac{\partial u}{\partial \eta} = \frac{1}{\eta} f(\tau).
\]
Integrăm în raport cu $ \eta $ și obținem, în final:
\[
  u(\tau, \eta) = f(\tau) \ln \eta + g(\tau) \Rightarrow %
  u(x, y) = f(xy) \ln x + g(xy).
\]

\vspace{1cm}

În unele cazuri, poate fi necesară o discuție după $ x, y $ pentru tipul ecuației:

\textbf{Exemplu 3:} Fie ecuația:
\[
  y \frac{\partial^2 u}{\partial x^2} + (x + y) \frac{\partial^2 u}{\partial x\partial y} + %
  x \frac{\partial^2 u}{\partial y^2} = 0.
\]

Deoarece $ A = y, B = \frac{x + y}{2}, C = x $, avem
\[
  \delta = AC - B^2 = \frac{-(x - y)^2}{4}
\]
și studiem separat pentru:
\[
  D_1 = \{(x, y) \in \RR^2 \mid \delta < 0 \}, \quad %
  D_2 = \{(x, y) \in \RR^2 \mid \delta = 0 \},
\]
care corespund, respectiv, cazurilor: \emph{hiperbolic}, pentru $ y \neq x $ și
\emph{eliptic}, pentru $ y = x $.

Mai departe, ecuația se rezolvă cu metodele cunoscute, corespunzătoare celor
două cazuri.

%%%%%%%%%%%%%%%%%%%%%%%%%%%%%%%%%%%%%%%%%%%%%%%%%%%%%%%%%%%%%%%%%%%%%%



\section{Coarda infinită. Metoda lui d'Alembert}

\indent\indent Pornim de la ecuația coardei infinite, care constă în determinarea funcției
$ u(x, t) $, definită pentru $ x \in \RR $ și $ t \geq 0 $, soluție a ecuației
coardei vibrante:

\[
  a^2 \frac{\partial^2 u}{\partial x^2} = \frac{\partial^2 u}{\partial t^2}, %
  \forall x \in \RR, t > 0.
\]

Presupunem că avem condiții inițiale, astfel că problema devine o problemă Cauchy:
\[
  u(x, 0) = \varphi(x), \quad \frac{\partial u}{\partial t}(x, 0) = \psi(x), \forall x \in \RR,
\]
cu $ \varphi $ și $ \psi $ funcții date.

Cum ecuația este deja în forma canonică, asociem ecuația caracteristică:
\[
  a^2 dt^2 - dx^2 = 0 \Rightarrow \frac{dt}{dx} = \pm\frac{1}{a}.
\]
Rezultă că familiile de curbe caracteristice sînt:
\[
  \begin{cases}
    x - at &= c_1 \\
    x + at &= c_2
  \end{cases}.
\]

Facem schimbarea de variabile corespunzătoare:
\[
  \begin{cases}
    \tau &= x - at \\
    \eta &= x + at
  \end{cases}
\]
și rezultă ecuația în forma canonică, în noile variabile:
\[
  \frac{\partial^2 u}{\partial \tau \partial \eta} = 0.
\]
Putem să o rezolvăm astfel:
\[
  \frac{\partial}{\partial \eta} \Big( \frac{\partial u}{\partial \tau} \Big) = 0 \Rightarrow %
  \frac{\partial u}{\partial \tau} = f(\tau).
\]
Acum putem integra în raport cu $ \tau $ și găsim:
\[
  u(\tau, \eta) = \int f(\tau) d\tau + \theta_2(\eta) \Leftrightarrow %
  u(\tau, \eta) = \theta_1(\tau) + \theta_2(\eta).
\]
Revenind la variabilele $ x, t $, avem:
\[
  u(x, t) = \theta_1(x + at) + \theta_2(x - at)
\]
și, folosind condițiile inițiale, avem:
\[
  \begin{cases}
    \theta_1(x) + \theta_2(x) &= \varphi(x) \\
    a\theta'_1(x) - a\theta'_2(x) &= \psi(x)
  \end{cases}.
\]
Integrăm a doua ecuație în raport cu $ x $ și obținem:
\[
  \begin{cases}
    \theta_1(x) + \theta_2(x) &= \varphi(x) \\
    a\theta_1(x) - a\theta_2(x) = \int_0^x \psi(\alpha) d\alpha + c
  \end{cases}
\]
Adunăm egalitățile și găsim:
\[
  \theta_1(x) = \frac{\varphi(x)}{2} + \frac{1}{2a} \int_0^x \psi(\alpha) d\alpha + \frac{c}{2a},
\]
iar prin scădere, găsim:
\[
  \theta_2(x) = \frac{\varphi(x)}{2} - \frac{1}{2a} \int_0^x \psi(\alpha) d\alpha - \frac{c}{2a}.
\]
Revenind la variabilele inițiale, avem:
\[
  \begin{cases}
    \theta_1(x + at) &= \frac{\varphi(x + at)}{2} + %
    \frac{1}{2a} \int_0^{x + at} \psi(\alpha) d\alpha + \frac{c}{2a} \\
    \theta_2(x - at) &= \frac{\varphi(x - at)}{2} - %
    \frac{1}{2a} \int_0^{x - at} \psi(\alpha) d\alpha - \frac{c}{2a}
  \end{cases}.
\]
Putem asambla soluția finală în forma:
\begin{equation} \label{eq:dalembert}
  u(x, t) = \frac{1}{2} \big[ \varphi(x - at) + \varphi(x + at) \big] + %
  \frac{1}{2a} \int_{x - at}^{x + at} \psi(\alpha) d\alpha,
\end{equation}
care se numește \textbf{formula lui d'Alembert}.

\begin{remark}\label{rk:ex-unic-coarda}
  Soluția problemei Cauchy asociată coardei vibrante există și este unică.
\end{remark}

%%%%%%%%%%%%%%%%%%%%%%%%%%%%%%%%%%%%%%%%%%%%%%%%%%%%%%%%%%%%%%%%%%%%%%

\section{Coarda finită. Metoda separării variabilelor (*)}

Pentru cazul lungimii finite a unei coarde, se folosește o metodă care este
atribuită lui Fourier și utilizează dezvoltări în serie. Această metodă se
numește \emph{metoda separării variabilelor}.

Pornim cu o problemă Cauchy similară, doar că lungimea coardei este conținută
într-un interval finit. Căutăm, deci, funcția $ u(x, t) $, definită pentru
$ 0 \leq x \leq l $ și $ t \geq 0 $, care satisface următoarele condiții:
\[
  \begin{cases}
    a^2 \frac{\partial^2 u}{\partial x^2} &= \frac{\partial^2 u}{\partial t^2} \\
    u(x, 0) &= \varphi(x) \\
    \frac{\partial u}{\partial t}(x, 0) &= \psi(x) \\
    u(0, t) &= u(l, t) = 0 \text{ (condiții la limită)}
  \end{cases}.
\]

Metoda de separare a variabilelor constă în găsirea unui șir infinit de soluții
de formă particulară, iar apoi, cu ajutorul acestora, formăm o serie ai cărei
coeficienți se determină în ipoteza ca suma seriei să dea soluția problemei tratate.

Soluțiile particulare se caută în forma:
\[
  u(x, t) = X(x)T(t)
\]
și cerem să satisfacă condițiile la limită:
\[
  \begin{cases}
    u(0, t) &= X(0)T(t) = 0 \\
    u(l, t) &= X(l)T(t) = 0
  \end{cases}
\]
Rezultă că vrem $ X(0) = X(l) = 0 $. Altfel, am avea $ T(t) = 0 $, ceea ce ar conduce la
soluția banală $ u(x, t) = 0 $.

Înlocuind în ecuația inițială, avem:
\[
  XT^{\prime\prime} = a^2X^{\prime\prime} T \Rightarrow %
  \frac{1}{a^2} \frac{T^{\prime\prime}}{T} = \frac{X^{\prime\prime}}{X}.
\]
Să remarcăm că în membrul stîng, funcția depinde doar de variabila $ t $, iar în
membrul drept, doar de variabila $ x $. Așadar, egalitatea nu poate avea loc
decît dacă ambele funcții sînt egale cu o constantă. Pentru conveniență, o vom nota
cu $ -\lambda $. Obținem ecuațiile:
\[
  \begin{cases}
    X^{\prime\prime} + \lambda X &= 0 \\
    T^{\prime\prime} + a^2\lambda T &= 0
  \end{cases}
\]
Prima dintre aceste ecuații este liniară, de ordinul al doilea, cu coeficienți
constanți. Soluția se obține:
\begin{itemize}
\item Dacă $ \lambda < 0 $, atunci
  \[
    X(x) = c_1 e^{-\sqrt{\lambda}x} + c_2 e^{-\sqrt{-\lambda} x},
  \]
  iar ținînd seama de condițiile la limită, avem:
  \[
    \begin{cases}
      c_1 + c_2 &= 0 \\
      c_1 e^{\sqrt{-\lambda} l} + c_2e^{-\sqrt{-\lambda} l} &= 0
    \end{cases},
  \]
  care se scrie echivalent:
  \[
    \begin{cases}
      c_1 + c_2 &= 0 \\
      c_1 e^{2\sqrt{-\lambda} l} + c_2 &= 0
    \end{cases}
  \]
  Determinantul matricei sistemului este nenul, deci el admite doar soluția banală.
\item Dacă $ \lambda = 0 $, atunci $ X(x) = c_1 x + c_2 $ și, ținînd seama de
  condițiile la limită, obținem din nou soluția banală.
\item Dacă $ \lambda > 0 $, soluția generală se scrie:
  \[
    X(x) = c_1 \cos \sqrt{\lambda} x + c_2 \sin \sqrt{\lambda} x.
  \]
  Din condițiile la limită, găsim:
  \[
    \begin{cases}
      X(0) &= c_1 = 0 \\
      X(l) &= c_2 \sin \sqrt{\lambda} l = 0.
    \end{cases}
  \]
  Din a doua ecuație, deducem că $ c_2 = 0 $, care conduce la soluția banală, sau
  $ \sin \sqrt{\lambda} l = 0 $, care înseamnă $ \lambda = \big(\frac{n\pi}{l}\big)^2 $.
  Putem scrie, atunci, soluția corespunzătoare acestei serii de valori în forma:
  \[
    X_n = c_n \sin \frac{n\pi}{l} x, \quad c_n \in \RR.
  \]
\end{itemize}

Înlocuim și integrăm acum ecuația după $ t $:
\[
  T^{\prime\prime} + a^2\Big( \frac{n\pi}{l} \Big)^2 T = 0,
\]
care are soluția generală:
\[
  T_n(t) = \alpha_n \cos \frac{n\pi}{l} at + \beta_n\sin \frac{n\pi}{l} at.
\]

Punînd laolaltă soluția după $ x $ și pe cea după $ t $, obținem:
\begin{equation} \label{eq:coarda-particulara}
  u_n(x, t) = X_n(x) T_n(t) = \Big( a_n \cos \frac{n\pi}{l} at + %
  b_n \sin \frac{n\pi}{l} at \Big) \cdot \sin \frac{n\pi}{l} x,
\end{equation}
unde $ a_n, b_n, c_n $ sînt constante ce provin din $ \alpha_n, \beta_n, c_n $.

Pentru a doua etapă a soluției, considerăm seria $ \sum u_n(x, t) $, adică:
\[
  \sum_{n = 1}^\infty \Big(a_n\cos\frac{n\pi}{l} at + %
  b_n \sin \frac{n\pi}{l} at\Big) \cdot \sin \frac{n \pi}{l} x.
\]

Presupunem că există $ u(x, t) $ suma seriei de mai sus, care este și soluția
problemei Cauchy, deci satisface și condițiile la limită, adică:
\[
  \begin{cases}
    u(x, 0) &= \sum_{n\geq 1} a_n \sin \frac{n\pi}{l} x = \varphi(x) \\
    \frac{\partial u}{\partial t}(x, 0) &= \frac{n\pi}{l}ab_n\sin\frac{n\pi}{l}x = \psi(x)
  \end{cases}.
\]

Putem privi aceste egalități ca dezvoltarea funcțiilor $ \varphi $ și $ \psi $ în serie
Fourier de sinusuri. Rezultă că putem afla coeficienții:
\begin{eqnarray} \label{eq:fourier}
  a_n &= \frac{2}{l} \int_0^l \varphi(x) \sin \frac{n\pi}{l} xdx \\
  b_n &= \frac{2}{n\pi a} \int_0^l \psi(x) \sin \frac{n \pi}{l} xdx.
\end{eqnarray}

\begin{remark}\label{rk:coarda-finita-fourier}
  Calculele de mai sus, împreună cu rezultate din teoria seriilor Fourier, ne asigură
  că funcția $ u(x, t) $ găsită este soluția problemei Cauchy.
\end{remark}

%%%%%%%%%%%%%%%%%%%%%%%%%%%%%%%%%%%%%%%%%%%%%%%%%%%%%%%%%%%%%%%%%%%%%%

\section{Exerciții}

1. Determinați soluția ecuației:
\[
  \frac{\partial^2 u}{\partial x^2} + 2 \frac{\partial^2 u}{\partial x \partial y} - %
  3 \frac{\partial^2 u}{\partial y^2} = 0,
\]
care satisface condițiile:
\[
  \begin{cases}
    u(x, 0) &= 3x^2 \\
    \frac{\partial u}{\partial y}(x, 0) &= \cos x
  \end{cases}.
\]

\emph{Soluție:} Cum $ AC - B^2 = -4 < 0 $, ecuația este de tip hiperbolic.
Din ecuația caracteristică, obținem schimbarea de variabile:
\[
  \begin{cases}
    \tau &= -3x + y \\
    \eta &= x + y
  \end{cases},
\]
iar forma canonică este $ \frac{\partial^2 u}{\partial \tau\partial \eta} = 0 $,
care are soluția:
\[
  u(\tau, \eta) = f(\tau) + g(\eta),
\]
cu $ f, g $ funcții de clasă $ \kal{C}^2 $, arbitrare. Revenind la variabilele
inițiale, avem:
\[
  u(x, y) = f(-3x + y) + g(x + y).
\]
Ținînd seama de condițiile inițiale din problema Cauchy, obținem sistemul:
\[
  \begin{cases}
    f(-3x) + g(x) &= 3x^2 \\
    f'(-3x) + g'(x) &= \cos x
  \end{cases}
\]
Integrăm a doua ecuație și avem:
\[
  \begin{cases}
    f(-3x) + g(x) &= 3x^2 \\
    -\frac{1}{3}f(-3x) + g(x) &= \sin x + c
  \end{cases}
\]
și prin schimbarea semnului primei ecuații și adunîndu-le, obținem:
\[
  \begin{cases}
    f(-3x) &= \frac{9}{4} x^2 - \frac{3}{4} \sin x - \frac{3c}{4} \\
    g(x) &= \frac{3}{4} x^2 + \frac{3}{4} \sin x + \frac{3c}{4}
  \end{cases}.
\]
Dacă notăm $ -3x = t $, atunci $ x = -\frac{t}{3} $ și găsim:
\[
  f(t) = \frac{t^2}{4} + \frac{3}{4} \sin \frac{t}{3} - \frac{3c}{4}.
\]
Așadar, soluția finală este:
\[
  \begin{cases}
    f(-3x + y) &= \frac{1}{4}(-3x + y)^2 + \frac{3}{4} \sin \frac{-3x + y}{3} - \frac{3c}{4} \\
    g(x + y)   &= \frac{3}{4}(x + y)^2 + \frac{3}{4} \sin(x + y) + \frac{3c}{4}
  \end{cases}.
\]
Rezultă că soluția problemei Cauchy este:
\[
  u(x, y) = \frac{1}{4}(-3x + y)^2 + \frac{3}{4} \sin \frac{-3x + y}{3} + %
  \frac{3}{4}(x + y)^2 + \frac{3}{4} \sin(x + y).
\]


\vspace{1cm}

2. Rezolvați ecuația:
\[
  3 \frac{\partial^2 u}{\partial x^2} + 7 \frac{\partial^2 u}{\partial x\partial y} + %
  2 \frac{\partial^2 u}{\partial y^2} = 0,
\]
cu condițiile:
\[
  \begin{cases}
    u(x, 0) &= x^3 \\
    \frac{\partial u}{\partial y}(x, 0) &= 2x^2
  \end{cases}.
\]

\vspace{1cm}

\emph{Soluție:} Ecuația este de tip hiperbolic, iar schimbarea de variabilă este:
\[
  \begin{cases}
    \tau &= 2x - y \\
    \eta &= x - 3y
  \end{cases},
\]
care conduce la forma canonică $ \frac{\partial^2 u}{\partial \tau\partial \eta} = 0 $,
de unde rezultă soluția generală:
\[
  u(x, y) = \varphi(2x - y) + \psi(x - 3y).
\]
Din condițiile problemei Cauchy, obținem:
\[
  \begin{cases}
    \varphi(2x) + \psi(x) &= x^3 \\
    -\varphi'(2x) - 3 \psi'(x) &= 2x^2
  \end{cases}.
\]
Integrăm a doua relație și obținem:
\[
  -\frac{1}{2} \varphi(2x) - 3 \psi(x) = \frac{2}{3}x^3 + k.
\]
Atunci:
\[
  \begin{cases}
    \varphi(2x) &= \frac{19}{96}(2x)^3 + c_1 \\
    \psi(x) &= -\frac{7}{12}x^3 - c_1
  \end{cases},
\]
de unde rezultă că soluția problemei Cauchy este:
\[
  u(x, y) = \frac{19}{96}(2x - y)^3 - \frac{7}{12}(x - 3y)^3.
\]

3. Rezolvați ecuația coardei vibrante infinite:
\[
  \frac{\partial^2 u}{\partial x^2} - \frac{\partial^2 u}{\partial t^2} = 0,
\]
cu condițiile inițiale:
\[
  \begin{cases}
    u(x, 0) &= \frac{x}{1 + x^2} \\
    \frac{\partial u}{\partial t}(x, 0) &= \sin x
  \end{cases}.
\]

\vspace{1cm}

\emph{Soluție:} Putem aplica direct formula lui d'Alembert (\ref{eq:dalembert}):
\begin{align*}
  u(x, t) &= \frac{1}{2} \Big[ \frac{x - t}{1 + (x - t)^2} + \frac{x + t}{1 + (x + t)^2} \Big] + %
            \frac{1}{2} \int_{x - t}^{x + t} \sin y dy \\
          &= \frac{1}{2} \Big[ \frac{x - t}{1 + (x - t)^2} + \frac{x + t}{1 + (x + t)^2} \Big] - %
            \frac{1}{2}[\cos(x + t) - \cos(x - t)] \\
          &= \frac{1}{2} \Big[ \frac{x - t}{1 + (x - t)^2} + \frac{x + t}{1 + (x + t)^2} \Big] - %
            \frac{1}{2} \Big[ -2\sin \frac{x + t + x - t}{2} \sin \frac{x + t - x + t}{2} \Big] \\
          &= \Big[ \frac{x - t}{1 + (x - t)^2} + \frac{x + t}{1 + (x + t)^2} \Big] + \sin x \sin t.
\end{align*}

\vspace{1cm}

4(*). Determinați vibrațiile unei coarde de lungime $ l $, avînd capetele fixate,
dacă forma inițială a coardei este dată de funcția:
\[
  \varphi(x) = 4 \Big( x - \frac{x^2}{l} \Big),
\]
iar viteza inițială este $ 0 $.

\vspace{1cm}

\emph{Soluție:} Aplicînd direct formula pentru coeficienții Fourier (\ref{eq:fourier}),
avem $ b_n = 0 $, iar
\[
  a_n = \frac{2}{l} \int_0^l 4\Big( x - \frac{x^2}{l} \Big) \sin \frac{n\pi}{l} xdx = %
  \frac{8}{l} \int_0^l x \sin \frac{n\pi}{l} xdx - %
  \frac{8}{l^2} \int_0^l x^2 \sin \frac{n\pi}{l}xdx.
\]

Calculăm:
\begin{align*}
  \int_0^l x\sin \frac{n\pi}{l} xdx &= -\frac{l}{n\pi} x \cos \frac{n\pi}{l}x \Big|_0^l + %
                                      \frac{l}{n\pi} \int_0^l \cos \frac{n\pi}{l} xdx \\
                                    &= -\frac{l^2}{n\pi}(-1)^n + %
                                      \frac{l^2}{n^2 \pi^2} \sin\frac{n\pi}{l} x\Big|_0^l \\
                                    &= (-1)^{n + 1} \frac{l^2}{n\pi}. \\
  \int_0^l x^2 \sin \frac{n\pi}{l} xdx &= -\frac{l}{n\pi}x^2 \cos \frac{n\pi}{l}x \Big|_0^l + %
                                         \frac{2l}{n\pi}\int_0^l x\cos \frac{n\pi}{l} xdx \\
                                    &= (-1)^{n + 1} \frac{l^3}{n\pi} + %
                                      \frac{2l}{n\pi}\Big[ \frac{l}{n\pi} x\sin \frac{n\pi}{l}x\Big|_0^l - \frac{l}{n\pi}\int_0^l \sin \frac{n\pi}{l}xdx \Big] \\
                                    &= (-1)^{n + 1}\frac{l^3}{n\pi} + %
                                      \frac{2l}{n^3 \pi^3}\cos \frac{n\pi}{l}x\Big|_0^l \\
                                    &= (-1)^{n + 1} \frac{l^3}{n\pi} + \frac{2l}{n^3 \pi^3}[(-1)^n + 1].
\end{align*}

Așadar:
\[
  a_n = (-1)^{n + 1} \frac{8l}{n\pi} - (-1)^{n + 1}\frac{8l}{n\pi} - %
  \frac{16}{n^3 \pi^3}[(-1)^n - 1],
\]
de unde obținem $ a_{2n} = 0, a_{2n+1} = \frac{32l}{(2n + 1)^3 \pi^3} $.

Punem laolaltă coeficienții și obținem soluția:
\[
  u(x, t) = \frac{32l}{\pi^3} \sum_{n \geq 0} \frac{1}{(2n + 1)^3} \cos\frac{(2n + 1)\pi}{l} t %
  \cdot \sin\frac{(2n + 1)\pi}{l} x.
\]

\vspace{1cm}

5(*). Rezolvați problema Cauchy asupra coardei vibrante finite:
\[
  \frac{\partial^2 u}{\partial t^2} = 4 \cdot \frac{\partial^2 u}{\partial x^2}, %
  \quad 0 < x < \pi, t > 0,
\]
cu condițiile inițiale și la limită:
\[
  \begin{cases}
    u(x, 0) &= \sin 3x - 4 \sin 10x \\
    \frac{\partial u}{\partial t}(x, 0) &= 2 \sin 4x + \sin 6x, 0 \leq x \leq \pi \\
    u(0, t) &= u(\pi, t) = 0, t \geq 0
  \end{cases}.
\]

\vspace{1cm}

\emph{Soluție:} Determinăm coeficienții din seria Fourier:
\[
  u(x, 0) = \sin 3x - 4\sin 10x \Rightarrow \sum a_n \sin nx = \sin 3x - 4 \sin 10x.
\]
Egalînd coeficienții, obținem $ a_3 = 1, a_{10} = -4, a_n = 0 $ în rest.

Mai departe:
\[
  \frac{\partial u}{\partial t}(x, 0) = 2\sin 4x + \sin 6x \Rightarrow %
  \sum 2nb_n \sin nx = 2\sin 4x + \sin 6x.
\]
Egalînd coeficienții, avem: $ b_4 = \frac{1}{4}, b_6 = \frac{1}{12}, b_n = 0 $ în rest.

Rezultă:
\[
  u(x, t) = \cos 6t \sin 3x - 4\cos 20t \sin 10x + \frac{1}{4} \sin 8t \sin 4x + %
  \frac{1}{12} \sin 12t \sin 6x.
\]

\vspace{1cm}

6(*). Aceeași cerință pentru:
\begin{enumerate}[(a)]
\item
  \[
    \frac{\partial^2 u}{\partial x^2} = \frac{1}{9}\frac{\partial^2 u}{\partial t^2}, %
    \quad 0 \leq x \leq \pi, t \geq 0,
  \]
  cu condițiile inițiale și la limită:
  \[
    \begin{cases}
      u(x, 0) &= \sin x \\
      \frac{\partial u}{\partial t}(x, 0) &= \sin x \\
      u(0, t) &= u(\pi, t) = 0
    \end{cases}.
  \]
\item
  \[
    \frac{\partial^2 u}{\partial x^2} - \frac{1}{4} \frac{\partial^2 u}{\partial t^2} = 0, %
    \quad 0 \leq x \leq 1, t \geq 0
  \]
  cu condițiile inițiale și la limită:
  \[
    \begin{cases}
      u(x, 0) &= x(1 - x) \\
      \frac{\partial u}{\partial t}(x, 0) &= 0 \\
      u(0, t) &= u(1, t) = 0
    \end{cases}.
  \]
\item
  \[
    \frac{\partial^2 u}{\partial x^2} - \frac{1}{4} \frac{\partial^2 u}{\partial t^2} = 0, %
    0 \leq x \leq \pi, t \geq 0
  \]
  cu condițiile inițiale și la limită:
  \[
    \begin{cases}
      u(x, 0) &= \sin 3x - 4\sin 10x \\
      \frac{\partial u}{\partial t}(x, 0) &= 2 \sin 4x + \sin 6x \\
      u(0, t) &= u(\pi, t) = 0
    \end{cases}.
  \]
\end{enumerate}

\end{document}